\documentclass[10pt,a4paper]{amsart}
\usepackage[margin=19mm]{geometry}
\usepackage{amsmath,amssymb,mathtools}
\usepackage[scr=rsfs,cal=euler]{mathalfa}
\usepackage{xfrac}
\usepackage[unicode,hidelinks,bookmarks=false]{hyperref}
\newcommand{\ie}{i.e.\ }
\newcommand{\cf}{cf.\ }
\newcommand{\resp}{respectively\ }
\newcommand{\Subsection}{\subsection}
\providecommand{\nobreakdash}{\nobreak\hbox{-}}
\setlength{\emergencystretch}{2em}
\multlinegap0pt
\usepackage{bm}
\usepackage{bbm}
\usepackage{dsfont}
\usepackage{xspace}
\usepackage{cancel}
\usepackage{mathtools}
\usepackage{amsthm}
\makeatletter
\@ifundefined{theorem}{\newtheorem{theorem}{Theorem}}{}
\@ifundefined{proposition}{\newtheorem{proposition}[theorem]{Proposition}}{}
\makeatother
\title[Hilbert's tenth problem for systems of diagonal quadratic fo]{Hilbert's tenth problem for systems of diagonal quadratic forms, and B\"{u}chi's problem}
\author{Stanley Yao Xiao}
\date{Source: arXiv:2412.16740v3 (CC BY 4.0)}
\begin{document}
\maketitle

\noindent\emph{Source and licence.} Hilbert's tenth problem for systems of diagonal quadratic forms, and B\"{u}chi's problem. Version arXiv:2412.16740v3 (CC BY 4.0), distributed under the Creative Commons Attribution 4.0 International licence, as recorded in the arXiv licence metadata for this exact version and shown on the abstract page https://arxiv.org/abs/2412.16740v3. Full rights notice: licenses/arxiv-2412.16740-CC-BY-4.0.txt.\par\smallskip

\noindent\textit{Editorial scope.} Counted unit: the Proposition whose source label is \texttt{None} in the article (source0.tex lines 406--407, source label \texttt{None}), delivered together with the article's own complete original proof of it (source0.tex lines 409--427, one \texttt{proof} environment of 19 lines). No mathematics is written, repaired or re-derived: statement and proof are the article's own text line for line. Required context retained uncounted: none. Every definition, display and label of those lines is retained unchanged. Omitted from this package and not redistributed: 1--405, 408--408, 428--1568. Nothing inside a retained excerpt is omitted or altered: every retained body is delivered exactly as the article's lines are written, so no omission receipt of any kind accompanies this package. Citations: the retained excerpts carry no citation command at all, so no bibliography entry of the article is transcribed and no reference is redistributed. Reference adaptation: none was needed and none was made; every reference inside the delivered unit resolves inside it, so no unresolved reference or citation is produced. Printed slips kept literal: no printed word, symbol, hypothesis or number of the counted statement or of its proof was altered. Layout and class: the article's own class is \texttt{amsart} with packages including amssymb, amsxtra, amsthm, amstext, amscd, amsfonts, fancyhdr, hyperref, enumerate, xcolor, comment, mathtools, moresize, amsmath; the wrapper is the lane's pinned \texttt{amsart} editorial wrapper at 10pt on A4, which restates the theorem-like environments the retained text needs and adds the article's own macro definitions that the retained text needs (none), with extra wrapper packages (bm, bbm, dsfont, xspace, cancel, mathtools). The delivered numbering is the unit's own. Assets: no retained span contains a figure environment, a graphics-inclusion command, a PDF-page inclusion, a table or a TikZ picture; no image file is copied into this package, the package declares no asset dependency and no image file is redistributed with it. Licence: version arXiv:2412.16740v3 of this article is distributed under the Creative Commons Attribution 4.0 International licence, as recorded in the arXiv licence metadata for this exact version and shown on the abstract page https://arxiv.org/abs/2412.16740v3. A full rights notice is delivered at licenses/arxiv-2412.16740-CC-BY-4.0.txt. Characters: every retained excerpt is pure ASCII, so the delivered engine, XeLaTeX with its default Latin Modern text font, reports no missing character.
\begin{proposition} Suppose $u_1, \cdots, u_n$ is a B\"{u}chi $n$-tuple. Then $u_1^2, \cdots, u_n^2$ is a non-singular monic quadratic progression of positive discriminant divisible by $16$. 
\end{proposition} 

\begin{proof}
	
	We write out the B\"{u}chi system
	\begin{equation} \label{orgsys} \begin{matrix} u_3^2 & - & 2 u_2^2 & + & u_1^2 & = & 2 \\ \vdots & & \vdots & & \ddots &  &  \vdots \\ x_{n}^2 & - & 2 u_{n-1}^2 & + & u_{n-2}^2 & = & 2.\end{matrix} 
	\end{equation}
	Note that (\ref{orgsys}), without the condition that the terms are square integers, defines a linear recurrence sequence $(a_n)_{n \geq 1}$ of order $2$ with characteristic polynomial $(t-1)^2$. Thus the general solution is given by 
	\[a_n = n^2 + \alpha n + \beta\]
	with $\alpha, \beta$ determined by initial conditions. Since we demand that $a_2 = u_2^2, a_1 = u_1^2$, we obtain the expression of $a_k = u_k^2$ as a quadratic polynomial: 
	\begin{equation} \label{quadpoly} u_k^2 = k(k-1) + k(u_2^2 - u_1^2) + u_1^2 = k^2 + k(u_2^2 - u_1^2 - 1) + u_1^2.
	\end{equation}
	For fixed $u_2, u_1$, the right hand side is a quadratic polynomial in $k$. It is clear that $u_2 \equiv u_1 + 1 \pmod{2}$, whence $2a = u_2^2 - u_1^2 - 1$ is even. We can then write the right hand side as 
	\[f(k) = k^2 + 2ak + b,\]
	where $b = u_1^2$. Then completing the square we have 
	\[f(k) = (k+a)^2 + b - a^2.\]
	The equation $u_k^2 = f(k)$ can then be expressed as 
	\[(u_k - k - a)(u_k + k + a) = b - a^2.\]
	
	
	In fact, we have that in the relevant case that $b - a^2 \equiv 0 \pmod{4}$. Indeed, recall that from (\ref{quadpoly}) that $b = u_1^2$ and $a = (u_2^2 - u_1^2 - 1)/2$. Clearly, $u_2, u_1$ have opposite parity. If $u_1$ is odd and $u_2$ is even, then $u_2^2 - u_1^2 - 1 \equiv 2 \pmod{4}$, hence $(u_2^2 - u_1^2 - 1)/2 = a$ is odd and hence $b - a^2$ is divisible by $4$. If $u_1$ is even and $u_2$ is odd, then $u_2^2 - u_1^2 - 1 \equiv 0 \pmod{4}$, and we again conclude that $b - a^2$ is divisible by $4$. \end{proof}

\end{document}
